\section[Atividades Desempenhadas Pelo Aluno no Projeto]{Atividades Desempenhadas Pelo Aluno no Projeto}

\subsection{Sprint 0}

Na Sprint 0 era previsto que o grupo realizasse os mockups do Figma, as
User Stories e preparar toda a ambientação, populando wiki e criando fontes de
comunicação, entregando uma prévia da primeira versão do trabalho, além de
preparar tudo para começar o desenvolvimento. Quanto a minha parte fiquei de
ajudar nos mockups e estudar, principalmente através da participação em
workshops. 

O grupo foi capaz de concluir as tarefas propostas e eu também
consegui concluir as minhas. Entrando em mais detalhes, estudei um pouco de
WebDevelopment(HTML e CSS básico), além de relembrar como funcionava o uso
básico do FIGMA, por fim em cooperação com um colega de equipe AGES I, projetei
8 telas(6 que foram para o MediumFi e outras 2 que eram apenas alternativas),
sendo 2 de criação de pergunta para RH(checkbox e múltipla escolha), 2 de
visualização de pergunta para o Usuário(as mesmas) e upload de conteúdo externo
para RH + visualização de conteúdo externo para o Usuário.

Durante a realização dessas tarefas os problemas foram os que mais eram
esperados, afinal padronizações e burocracias são problemas bem comuns para
quem está trabalhando pela primeira vez em um projeto grande e uma equipe
grande, e assim foi para mim, deixar tudo padronizado e até ter medo de fazer a
primeira ação por não ter certeza de como o time trabalharia foram os maiores
problemas.

Por outro lado, achei interessante aprender mais a trabalhar em time e como
a criação de componentes/fontes/cores padronizados são essenciais para manter o
projeto bem estruturado e que usá-las ao próprio favor consegue poupar muito
tempo e deixar muito mais interligado o projeto.

Seguindo os aprendizados pretendo manter esta lógica de utilização de
componentes, reaproveitando e padronizando o máximo de código de meus outros
companheiros no FrontEnd para poder manter o estilo de programação uniforme,
ainda mais se utilizando dos plugins Prettier e ESLint, adotados pelo time para manter a uniformidade do código.
O começo sendo um AGES I não participei nas decisões da equipe, mas ao saber que teríamos que apresentar Mockups e apresentar para a cliente um modelo base de como estaríamos desenvolvendo a aplicação, focando em sua aparência no início, busquei aprender sobre a ferramenta Figma, onde poderia ser desenvolvido o design da aplicação e o protótipo dele, para que fosse possível apresentar para a cliente.

Na fase inicial, havia muitos tipos de protótipos feitos por diversas pessoas, como o ambiente não foi totalmente organizado e não havia equipes definidas, ao final da sprint metade das telas foram descartadas e a cliente escolheu qual delas iríamos manter.

Durante o desenvolvimento dos mockups, os AGES III estavam decidindo as tecnologias que iriamos usar para desenvolver o back-end do aplicativo, onde foi apresentada a proposta de escolha da linguagem de programação que utilizaríamos.

Minha decisão foi utilizarmos Python pois por mais que todos os AGES I estavam familiarizados com a linguagem Java, os nossos AGES III possuíam uma afinidade maior com Python, sabendo disso, propus ela, como uma chance de aprender uma nova linguagem e facilitaria muito mais receber ajuda deles em quando eles sabem mais sobre o assunto.

	Ao final da Sprint, graças ao resto do time, eu não tive grandes dificuldades em me integrar aos meus futuros estudos das tecnologias que iriamos passar a utilizar para a realização das tarefas, assim eu me preparei para estudar sobre Fast API e Python, que eram focadas para o back-end, algo que eu possuía interesse em aprender sobre.
\subsection{\textit{Sprint} 1}

Para esta \textit{sprint}, as atividades previstas para a equipe eram fazer uma tela
de \textit{login} e uma tela \textit{home}. Quanto às atividades previstas para meu \textit{squad}, estava
prevista a implementação do \textit{front-end} da tela de \textit{login}, utilizando React
\cite{react-docs} e TypeScript \cite{typescript-docs},
utilizando componentes genéricos criados pelo grupo, como um botão, por
exemplo, e em alguns casos até criando componentes genéricos, como \textit{inputs}.

A equipe inteira conseguiu completar as \textit{tasks} e a parte do \textit{back-end}
conseguiu avançar até mais do que o previsto. Quanto ao meu \textit{squad}, conseguimos
implementar a tela de \textit{login} quase igual ao protótipo do Figma \cite{figma-help} e criar um novo
componente genérico de \textit{input}.

Por incrível que pareça não foram encontrados tantos problemas, creio que o
fato de a tela desenvolvida ser simples facilitou o desenvolvimento sem muitas
barreiras, apenas com algumas correções para garantir que a página fosse
responsiva e o mais genérica possível.

Quanto às lições, aprendi que a maioria dos itens utilizados na página devem
ser transformados em componentes com o intuito de facilitar o trabalho dos outros
colegas e até meu trabalho no futuro.

Prosseguirei, então, com minha equipe, que está planejando fazer cerca de 8
telas no \textit{front-end} e todo o \textit{back-end} por trás das que ainda não foram feitas (parte
destas 8 telas já tem o \textit{back-end} pronto). Quanto ao meu planejamento, estou
esperando meu \textit{squad} receber uma \textit{task}; no entanto, creio que serão cerca de 3 a 4 telas para meu \textit{squad} desenvolver.

Na Sprint 1 eu comecei a realizar meus estudos onde eu foquei mais em Python. Nessa Sprint foi realizado a definição das equipes, sendo três equipes para cada AGES III, onde a primeira equipe iria dar mais importância para o front-end, a segunda para o back-end e a terceira seria uma mistura dos outros dois. Sabendo disso eu decidi escolher a equipe do back-end e continuar meus estudos.

	Dediquei um pouco do meu tempo para sanar dúvidas com o líder da minha equipe em relação ao estabelecimento do ambiente que estaríamos utilizando para trabalhar no projeto, onde após o estabelecimento foi disponibilizado as tarefas que seriam feitas, e na primeira sprint por ser no começo, havia pouquíssimas tarefas para o back-end, sendo eu fiquei encarregado de tarefa de: “Validação de Dados do Cadastro”.
 
	Tive um pouco de dificuldade em realizar minha tarefa pois eu tive que utilizar método de validação conhecido como Regex, algo que eu não possuía familiaridade, felizmente recebi ajuda de membros da minha equipe para então conseguir utilizar esse método. Entretanto ao realizar o envio da minha tarefa para análise foi encontrado erros que giravam em volta do Regex e da minha falta de atenção, o que acarretou a demora para a segunda análise da minha tarefa, que não foi finalizada a tempo de a Sprint acabar. Isso me fez ter um pouco mais de atenção à comunicação e ao cuidado com o tempo para a entrega da tarefa. Entretanto tive experiências positivas em relação a linguagem de Python, que despertou meu interesse em aprender mais sobre ela.
\subsection{\textit{Sprint} 2}

Para a \textit{sprint} 2, o time havia previsto a criação de cerca de 6 telas com seu
\textit{back-end} respectivo, entre elas criação e gerenciamento de módulos, setores,
cursos, atividades e colaboradores, além da tela inicial. Além disso, haviam sido
previstas 2 atividades de \textit{back-end}: autenticação e envio de \textit{e-mail}. Quanto ao meu time,
havia sido previsto o desenvolvimento de um componente \textit{table}, além da página de
administrar os colaboradores.

O time desempenhou muito bem em relação às expectativas, visto o tanto de
problemas encontrados no desenvolvimento. Das 6 páginas propostas é possível
dizer que mais de 95\% do prometido foi entregue, faltando apenas detalhes de
padronização de componentes e estilização de páginas, ficando 2 \textit{front-ends} e 1
\textit{back-end} com previsão de mudanças. Quanto ao meu time, aprendemos o
componente \textit{table} e seu funcionamento, além de implementá-lo como componente
geral no projeto e construir a página de administração com tabela e filtros. Além
disso, criamos o modal da página de administração. Por fim, foi possível começar a
pensar na integração para a página de administração.

O principal problema encontrado pelo time vai além da programação. Desde
o começo do mês de maio de 2024, o estado do Rio Grande do Sul, incluindo
principalmente cidades que fazem fronteira com Porto Alegre, enfrentaram uma das
maiores enchentes da história, paralisando as aulas por um tempo e preocupando
muitos. O tempo ocioso junto com a preocupação não foi tanto problema para esta
\textit{sprint}, já que, na hora em que isso aconteceu, a maioria dos itens da entrega já estava
pronta; no entanto, tende a retardar um pouco a velocidade de desenvolvimento do
início da próxima \textit{sprint}. Além disso, entregas de trabalhos e provas foram uma
preocupação paralela à da \textit{sprint} que inevitavelmente atrapalharam um pouco o
desenvolvimento no começo da \textit{sprint}, empurrando a maior parte do
desenvolvimento para a última metade.

Cada vez mais, conciliar a Ages com estudo e estágio tem se provado uma
tarefa que requer planejamento adiantado, principalmente em uma \textit{sprint} que teve
tantas incertezas como esta, a habilidade de retraçar o planejamento e se adaptar à
situação é colocada em prova, ganhando experiência e aprendendo a agir com
calma mediante situações complicadas e imprevistos.

A equipe planeja terminar os últimos ajustes que ficaram faltando nas páginas
pendentes da \textit{sprint} passada, além de continuar desenvolvendo as telas que faltam e
adiantar bem a parte da integração. Quanto à minha parte, ainda não foi designada
nenhuma tarefa de imediato, mas já tenho em mente continuar a parte da integração
da página de administração, além de colaborar com o time nas telas que forem
posteriormente designadas.

Na Sprint 2 começaram a ocorrer alguns problemas em relação à estrutura do banco de dados. Para começar, nosso banco de dados utilizava uma estrutura de muitos para muitos, isso significa que não existe uma tabela de listas de itens, e sim um amontoado de itens que correspondem a um evento específico de um usuário, isso dificultou nas minhas tarefas, já que para que eu pudesse retornar todas as listas de um usuário ou até mesmo deletar uma, sempre seria preciso o evento, seria o único meio de distinguir qual item pertence a uma lista ou outra.

	Nesse período ocorreu uma tragédia no Rio Grande do Sul, onde muitas pessoas foram afetadas pelas enchentes, o que fez com que o projeto ficasse parado por duas semanas. Esse tempo de inatividade causou um aperto no prazo de entrega da Sprint, e muitas tarefas foram deixadas para a próxima Sprint. Nesse tempo os meus Ages IV refizeram o planejamento do projeto e tivemos que sacrificar muitas funcionalidades do aplicativo para que fosse possível entregar ele funcionando.
 
	Durante esse intervalo, conversei com meu time sobre a melhor solução para minhas tarefas, uma das conclusões seria criar a tabela de lista de itens, o que facilitaria muitas outras tarefas, mas também teríamos que refazer algumas tarefas já concluídas. 
\subsection{\textit{Sprint} 3}

Para a penúltima \textit{sprint}, o time havia previsto a integração de todas as telas
do \textit{front-end} com o \textit{back-end}, além de entregar os débitos técnicos que haviam ficado
da \textit{sprint} passada. Quanto ao meu \textit{squad}, ficamos com a integração da página de
administração que já havia sido iniciada, além de assumir a criação dos débitos
técnicos dos modais de edição e deleção de usuários.

Quanto aos débitos técnicos, o time inteiro foi capaz de entregar.
Especificamente da minha parte, os modais de edição e deleção de usuário foram
entregues com certa facilidade. Já as integrações foram uma dor de cabeça para
fazer, sendo perceptível que nenhuma página chegou a ser completamente
integrada. Pelo menos, é possível dizer que todas as integrações planejadas estão
bem alinhadas e não devem ser muito complicadas de finalizar.

Apesar de acreditar que a inércia de código seria um dos maiores problemas,
ela não foi fator, pelo menos não para o meu \textit{squad}. No entanto, foi possível ver
membros do time desanimados e partes do projeto andando mais devagar do que
deveriam. Tanto minha equipe quanto as outras do \textit{front-end} encontraram certas
complicações na integração, em parte por inexperiência, mas também em parte por rotas
não implementadas do \textit{back-end}.

Como integrante do time, percebi que dependo dos outros integrantes do time
e do trabalho deles. Sendo assim, mediante momentos complicados de desmotivação
do time, aprendi que é necessário dar um passo à frente e se preocupar com o
desenvolvimento dos outros companheiros de time, zelando pelo ânimo dos
companheiros e tentando incentivá-los sempre que possível para garantir que o time
esteja trabalhando de maneira uniforme.

Além disso, a persistência se mostrou uma grande aliada quanto à
resolução de muitos \textit{bugs}. Da parte técnica, aprendi fontes valiosas para obter
informações do projeto, como por exemplo, a documentação do Swagger
\cite{swagger-docs}. Também
tive muito contato com a aba de desenvolvedor do navegador, já que foi necessário
depurar a parte das requisições, além de tratar diversos códigos de erro.

Com a \textit{sprint} final se aproximando, as \textit{tasks} que sobraram são as pendências,
além de dar uma refinada no código visualmente e estruturalmente.
Por enquanto, recebi a \textit{task} de continuar com a integração da página de
administrador, além de adicionar um botão para baixar o relatório e fazer este botão funcionar. Com
a finalização dessas \textit{tasks}, falarei com outros membros da equipe e trabalharei em alguma das \textit{tasks} pendentes.

A terceira Sprint foi uma das mais curtas do projeto inteiro, isso se deve por ela ter ocorrido durante as enchentes e também pelo fato que ela consistia de poucas tarefas novas e muitas das Sprint passada que já estavam no processo de conclusão. 

Eu tive dificuldades em testar o código no meu ambiente, mas felizmente meus colegas da minha equipe me ajudaram a validar muito desses códigos, em uma das minhas tentativas de fazer minha tarefa, eu tentei criar o modelo de listas no código, mas minha inexperiência fez eu ter muitos problemas e acabei não conseguindo concluir o código através desse caminho. Por fim eu decidi utilizar a ideia inicial de requisitar o evento para reunir os itens, mas com a falta de tempo a minha tarefa não foi aprovada nessa Sprint.
\subsection{Sprint 4}

Por se tratar da sprint final, não havia muitas tarefas novas, mas sim débitos
técnicos de sprints passadas, além de complicações com a integração que deveriam
ser resolvidas. Da minha parte tinha a página de administração, na qual eu
precisava atualizar parâmetros e outros atributos do display conforme redefinições
de regra de negócio e alterações do backend, além de adicionar qualquer botão que
faltasse na tela de administração.

O time foi capaz de entregar praticamente tudo, incluindo tarefas novas que
iam sendo criadas conforme reflexão do próprio time relativas à experiência de
usuário. Da minha parte fiquei um bom tempo esperando a rota do backend ser
implementada para a parte de administração, já que era uma das páginas mais
complexas da parte do backend.

Enquanto esperava o backend finalizar a rota, me prontifiquei para realizar pequenas mudanças de qualidade de vida para a página,
como adicionar um botão de relatório integrado com o backend, reunir com outros
colegas para adicionar o drag and drop e a criação de usuários por arquivo e outras
tarefas do tipo. Por fim participei frequentemente das reuniões de alinhamento do
time fora da aula para garantir que me manteria informado e sempre ajudando o time nas tarefas necessárias.

Más práticas por parte do backend como, não padronizar o nome esperado do
payload e deixar a documentação do swagger mal-feita, foram os principais
problemas encontrados nesta última sprint, já que para realizar as integrações era necessário
entender o código em NestJS \cite{nestjs-docs} ou perguntar detalhadamente para o integrante
responsável pelo código e isso foi responsável por diminuir um pouco a velocidade
com que o desenvolvimento ocorreu.

Da minha parte acredito que poderia ter ido atrás de aprender o mínimo de NestJS
\cite{nestjs-docs},
já que isso teria facilitado muito a legibilidade do código do backend e, portanto,
facilitado a parte de integração. Ao mesmo tempo pensando pelos olhos de quem
faz o backend, que é onde planejo me concentrar mais na próxima Ages, entendi
muito bem qual a importância de uma boa documentação e um Swagger fácil de
entender, já que percebi na prática o quanto ele fez falta para alguém do frontend
que não faz ideia de como funciona a linguagem do backend.

Com a última sprint desta Ages finalizada, só basta pensar na próxima, como
citado acima, espero me envolver muito mais no backend para chegar de Ages 3
com uma boa experiência nos dois lados. Além disso acredito que a função de Ages 2 de
cuidar do banco de dados combina muito bem com o desenvolvimento backend, já
que os dois referem-se bem restritamente à regra de negócio. Além disso estou
querendo aprender Java Spring \cite{spring-docs} por fora da aula, por interesse próprio, o que me
daria um bom conhecimento básico do backend.

A última Sprint foi a mais preocupante, onde todos focaram em concluir e máximo de funcionalidades possíveis antes de mostrar para o nosso cliente. Todos estavam bem empenhados em concluir suas tarefas, eu fiquei corrigindo os erros que meu líder de time detectava nas minhas tarefas.

Em uma das reuniões fora das aulas, foi decidido que seria melhor criar a tabela de itens, portanto todas as tarefas em relação ao back-end foram refeitas, muitas delas foi pelo AGES II do meu time. E com a etapa final de correção de bugs completa, conseguimos entregas um projeto bem satisfatório para nosso cliente.